Eine Stabhochspringerin sprintet mit \vAO; ihr Schwerpunkt ist \hAO\ über dem Boden. Der Stab wandelt ihre kinetische Energie in Lageenergie um (er biegt sich und gibt die gespeicherte Energie wieder ab).

\begin{abcliste}
\abc
Wie hoch kann ihr Schwerpunkt höchstens steigen?

\abc
Wie schnell müsste sie laufen, damit ihr Schwerpunkt bis auf \hZO\ steigt?
\end{abcliste}

Vernachlässige Reibung und Luftwiderstand. Rechne mit \gO.
